\documentclass[12pt,a4paper]{article}
\usepackage[margin=25mm, headheight=15pt, headsep=10pt]{geometry}
\usepackage{fontspec}
\usepackage[table]{xcolor} % Note: table option is required for rowcolors
\usepackage{titlesec}
\usepackage{tocloft}
\usepackage{fancyhdr}
\usepackage{booktabs}
\usepackage{tcolorbox}
\tcbuselibrary{skins, breakable}
\usepackage[colorlinks=true, linkcolor=emerald, urlcolor=emerald, bookmarks=true]{hyperref}

\definecolor{emerald}{RGB}{0,102,51}
\definecolor{darkred}{RGB}{139,0,0}
\definecolor{lightgray}{RGB}{245,245,245}

\usepackage[nil]{babel}
\babelprovide[import, main]{bengali}
\babelprovide[import]{arabic}
\babelfont{rm}[Renderer=HarfBuzz,Scale=1.0]{Noto Serif Bengali}
\babelfont{sf}[Renderer=HarfBuzz]{Noto Sans Bengali}
\babelfont{tt}[Renderer=HarfBuzz]{Noto Sans Bengali}
\babelfont[arabic]{rm}[Renderer=HarfBuzz,Scale=1.1]{Noto Naskh Arabic}

% Section styling
\titleformat{\section}{\Large\bfseries\color{emerald}}{\thesection.}{0.5em}{}
\titleformat{\subsection}{\large\bfseries\color{emerald!85!black}}{\thesubsection.}{0.5em}{}
\titleformat{\subsubsection}{\normalsize\bfseries\color{emerald!70!black}}{\thesubsubsection.}{0.5em}{}
\titlespacing*{\section}{0pt}{18pt}{8pt}

% Fancy Header/Footer setup
\pagestyle{fancy}
\fancyhf{} % clear all header and footer fields
\fancyhead[L]{\color{emerald}\small\textbf{\nouppercase{\leftmark}}}
\fancyfoot[L]{\scriptsize\color{black!50}{\lat{AI}}-সংকলিত, আলেম-যাচাইকৃত নয়}
\fancyfoot[C]{\color{emerald!70!black}\thepage}
\renewcommand{\headrulewidth}{0.4pt}
\renewcommand{\headrule}{\hbox to\headwidth{\color{emerald}\leaders\hrule height \headrulewidth\hfill}}
\renewcommand{\footrulewidth}{0pt}

% Custom boxes
% 1. Ayat/Hadith Box
\newtcolorbox{ayatbox}[1][]{
    enhanced, breakable, 
    colback=emerald!5, 
    colframe=emerald, 
    boxrule=0.5pt, 
    arc=3pt, 
    left=10pt, right=10pt, top=10pt, bottom=10pt,
    #1
}

% 2. Quote/Translation Box
\newtcolorbox{quotebox}[1][]{
    enhanced, breakable,
    frame hidden,
    colback=lightgray,
    borderline west={3pt}{0pt}{emerald},
    left=15pt, right=10pt, top=10pt, bottom=10pt,
    sharp corners,
    #1
}

% 3. AI-generated notice box
\newtcolorbox{warnbox}[1][]{
    enhanced, breakable,
    colback=red!4,
    colframe=darkred,
    boxrule=0.8pt,
    arc=3pt,
    left=10pt, right=10pt, top=10pt, bottom=10pt,
    #1
}

% Commands
\newcommand{\arab}[1]{\foreignlanguage{arabic}{#1}}
\newcommand{\kib}[1]{\textcolor{emerald}{\textbf{#1}}}

% Latin (English) fallback: Noto Serif Bengali has NO Latin glyphs
\newfontfamily\latfont[Renderer=HarfBuzz]{Noto Serif}
\newcommand{\lat}[1]{{\latfont #1}}

\setlength{\parskip}{5pt}
\setlength{\parindent}{0pt}

% Fix for missing bullet in Noto Serif Bengali
\renewcommand{\labelitemi}{$\bullet$}



\begin{document}
\begin{center}{\Large\bfseries\color{emerald} পার্ট ১ — সালাতের মর্যাদা ও গুরুত্ব}\end{center}
\vspace{1em}
\section*{পার্ট ১ — সালাতের মর্যাদা ও গুরুত্ব}

\begin{warnbox}
 \textbf{গুরুত্বপূর্ণ নোটিশ (\lat{Important} \lat{Notice}):} এই কন্টেন্টটি \lat{AI} (কৃত্রিম বুদ্ধিমত্তা) দিয়ে প্রস্তুত ও সংকলিত। \lat{AI} কঠোর সূত্র-মান নীতি মেনে শুধু নির্ভরযোগ্য উৎস থেকে তথ্য সংগ্রহ করেছে — কেবল সহীহ/হাসান হাদিস ও সহীহ সনদের আসার রাখা হয়েছে, দুর্বল সনদ বাদ দেওয়া হয়েছে। তবে এটি কোনো আলেম বা মুহাদ্দিস কর্তৃক যাচাইকৃত নয়।
\end{warnbox}


\begin{quote}
\textbf{পড়ার নিয়ম:} এই ফাইলটি পার্ট ১ — সালাতের মর্যাদা ও গুরুত্ব। আয়াতগুলোর পূর্ণ তাফসীরের জন্য দেখুন `কুরআন/\lat{01}-কুরআনের-আয়াত.\lat{md}`; প্রতিটি হাদিসের সূত্র কার্ডের জন্য দেখুন `হাদিস/\lat{02}-হাদিস-ও-আসার.\lat{md}`।
\end{quote}


\medskip\hrule\medskip

\subsection*{১. কুরআনে সালাতের মর্যাদা}

কুরআনে সালাত এমন একটি ইবাদত, যাকে আল্লাহ নিজের \textbf{স্মরণের (\lat{remembrance})} সঙ্গে যুক্ত করেছেন — \lat{“}নিশ্চয়ই আমিই আল্লাহ... আমার স্মরণের জন্য নামাজ কায়েম করো\lat{”} (সূরা ত্বহা ২০:১৪)। অর্থাৎ নামাজ শুধু বিধান নয়; এটি আল্লাহর সাথে বান্দার জীবন্ত সংযোগ।

\textbf{মূল আয়াতগুলো এক নজরে:}

\begin{table}[h]\centering\small
\begin{tabular}{ll}
আয়াত & বার্তা \\
\hline
ত্বহা ২০:১৪ & নামাজ — আল্লাহর স্মরণের মাধ্যম \\
বাকারাহ ২:৪৩ & নামাজ-জাকাত একসাথে; জামাতের নির্দেশ \\
নিসা ৪:১০৩ & নামাজ নির্দিষ্ট সময়ে \textbf{ফরজ (\lat{obligatory})} \\
আনকাবুত ২৯:৪৫ & নামাজ ফাহেশা-মুনকার থেকে বিরত রাখে \\
মু'মিনুন ২৩:১-২ & সফল মুমিনের প্রথম গুণ — নামাজে খুশু \\
বাকারাহ ২:২৩৮ & সব নামাজে যত্ন, বিশেষত আসর \\
মাউন ১০৭:৪-৫ & নামাজে উদাসীনতার ধিক্কার \\
\hline
\end{tabular}
\end{table}

এই আয়াতগুলো থেকে দুটি বিষয় স্পষ্ট: (১) নামাজ ইসলামের \textbf{কাঠামো (\lat{structure})} — যা ছাড়া দ্বীন অসম্পূর্ণ; (২) নামাজের মান (খুশু, সময়, পূর্ণতা) আল্লাহর কাছে সংখ্যার মতোই গুরুত্বপূর্ণ।

\medskip\hrule\medskip

\subsection*{২. সুন্নাহয় সালাতের মর্যাদা}

নবীজি (সাল্লাল্লাহু আলাইহি ওয়া সাল্লাম)-এর বাণীতে সালাতের মর্যাদা আরও বিস্তারিত:

\begin{itemize}
\item \textbf{প্রথম হিসাব:} \lat{“}কিয়ামতের দিন বান্দার আমলের মধ্যে সর্বপ্রথম নামাজের হিসাব নেওয়া হবে...\lat{”} (তিরমিযী ৪১৩; আবু দাউদ ৮৬৪; নাসাঈ ৪৬৫ — সহীহ)
\item \textbf{দ্বীনের স্তম্ভ:} \lat{“}দ্বীনের মাথা ইসলাম, তার স্তম্ভ নামাজ...\lat{”} (তিরমিযী ২৬১৬ — হাসান সহীহ)
\item \textbf{চোখের শীতলতা:} \lat{“}চোখের শীতলতা রাখা হয়েছে নামাজে\lat{”} (নাসাঈ ৩৯৩৯ — হাসান)
\item \textbf{শেষ অসিয়ত:} মৃত্যুশয্যায় — \lat{“}নামাজ, নামাজ! আর তোমাদের অধীনস্থদের ব্যাপারে আল্লাহকে ভয় করো\lat{”} (আবু দাউদ ৫১৫৬ — সহীহ)
\item \textbf{ফজরের দুই রাকাত:} \lat{“}দুনিয়া ও তার সবকিছুর চেয়ে উত্তম\lat{”} (মুসলিম ৭২৫)
\item \textbf{পাঁচ নামাজ কাফফারা:} \lat{“}পাঁচ ওয়াক্ত নামাজ ও জুমা... মধ্যবর্তী গুনাহের কাফফারা\lat{”} (মুসলিম ২৩৩)
\end{itemize}
মজার বিষয় — নবীজির জীবনে নামাজ কখনো \textbf{\lat{boring} (একঘেয়ে)} \lat{routine} ছিল না; বিপদে তিনি নামাজে দাঁড়াতেন (আবু দাউদ ১৩১৯), নামাজকে বলতেন চোখের শীতলতা। নামাজ তার জন্য বিশ্রাম ছিল — যেমনটা তিনি বিলালকে বলেছিলেন: \lat{“}হে বিলাল, আমাদের নামাজে প্রশান্তি দাও\lat{”} (আবু দাউদ ৪৯৮৫ — সহীহ)।

\medskip\hrule\medskip

\subsection*{৩. সালাত ও ঈমান}

সালাত ও ঈমানের সম্পর্ক এত গভীর যে, নবীজি (সা.) নামাজ ত্যাগকে ঈমানের সীমার সাথে যুক্ত করেছেন:

\begin{quote}
\textbf{\lat{“}এক ব্যক্তি ও শিরক-কুফরের মধ্যে পার্থক্য হলো নামাজ ত্যাগ করা।\lat{”}} — সহীহ মুসলিম ৮২; তিরমিযী ২৬২০
\end{quote}

\begin{quote}
\textbf{\lat{“}আমাদের এবং তাদের মধ্যে অঙ্গীকার হলো নামাজ; যে তা ত্যাগ করল, সে কুফরি করল।\lat{”}} — তিরমিযী ২৬২১; ইবনে মাজাহ ১০৭৯ (সহীহ)
\end{quote}

\textbf{এই হাদিসের ব্যাখ্যা নিয়ে আলেমদের তিনটি অবস্থান} (সৎভাবে তুলে ধরলাম):

\begin{enumerate}
\item \textbf{আক্ষরিক মত (সালাফের একাংশ):} নামাজ ত্যাগ করা নিজেই কুফর — ব্যক্তি দীন থেকে বের হয়ে যায়। ইবনে হাযম, ইমাম আহমাদের একটি প্রসিদ্ধ মত, শাইখুল ইসলাম ইবনে তাইমিয়্যার প্রসিদ্ধ মত ও তার ছাত্র ইবনুল কাইয়্যিমের মত — সতর্কতামূলকভাবে নামাজ ত্যাগকারীকে কাফির গণ্য করা।
\end{enumerate}
\begin{enumerate}
\item \textbf{জমহুর ফকীহদের মত (হানাফী, শাফি'ঈ, মালিকী, হাম্বলী মাযহাবের প্রসিদ্ধ অবস্থান):} যে ব্যক্তি অলসতা/অবহেলায় নামাজ ছাড়ে কিন্তু নামাজ ফরজ ও আবশ্যক মনে করে — সে \textbf{কাফির নয়}, বড় গুনাহগার মুসলিম; তাকে তাওবা ও পুনরায় নামাজের আহ্বান জানানো হবে।
\end{enumerate}
\begin{enumerate}
\item \textbf{মধ্যপন্থা (বহু সালাফ ও মুহাদ্দিস):} নামাজ ত্যাগ দুই প্রকার — (ক) নামাজ ফরজ মনে না করা বা অবজ্ঞা করা = কুফর, (খ) ফরজ মানলেও ছেড়ে দেওয়া = বড় গুনাহ, কুফর নয়। এই পার্থক্যটিই আধুনিক যুগের অধিকাংশ গবেষকের \textbf{\lat{preferred} (পছন্দের)} অবস্থান।
\end{enumerate}
\textbf{গুরুত্বপূর্ণ সতর্কতা:} তিনটি মতই নামাজ ত্যাগকে সবচেয়ে ভয়াবহ অন্যায় মনে করে। মতভেদটি \textbf{শাস্তির শ্রেণি} নিয়ে, নামাজের গুরুত্ব নিয়ে নয়। কোনো মতই নামাজ ত্যাগকে হালকা করার \textbf{\lat{license} (অনুমতি)} দেয় না।

\medskip\hrule\medskip

\subsection*{৪. সালাত ও দায়বদ্ধতা (\lat{Accountability})}

নামাজ কেবল মর্যাদার বিষয় নয় — এটি \textbf{হিসাবের (\lat{reckoning})} প্রথম দরজা:

\begin{itemize}
\item কিয়ামতে প্রথম জিজ্ঞাসা নামাজের — তিরমিযী ৪১৩ (সহীহ)
\item নামাজের ঘাটতি নফল দিয়ে পূরণের সুযোগ — একই হাদিস
\item আসর ছুটলে পরিবার-ধন হারানোর সাদৃশ্য — বুখারি ৫৫২, মুসলিম ৬২৬
\item মুনাফিকদের চিহ্ন — ইশা ও ফজরের জামাতে অনীহা — বুখারি ৬৫৭, মুসলিম ৬৫১
\end{itemize}
দায়বদ্ধতার দিক থেকে নামাজের তিনটি স্তর:

\begin{enumerate}
\item \textbf{ফরজ নামাজ} — যে ছাড়বে তার জবাবদিহি কঠিন; কবরে ও কিয়ামতে প্রথম প্রশ্ন।
\item \textbf{সুন্নত নামাজ} — নবীজির পথ; যা ফরজের ঘাটতি পূরণ করে (তিরমিযী ৪১৩-এর ব্যখ্যা)।
\item \textbf{নফল নামাজ} — আল্লাহর নৈকট্যের অতিরিক্ত \textbf{পুঁজি (\lat{capital})}; \lat{“}ফরজের পর সর্বোত্তম নামাজ রাতের নামাজ\lat{”} (মুসলিম ১১৬৩)।
\end{enumerate}
\medskip\hrule\medskip

\subsection*{৫. নামাজ ত্যাগের সতর্কতা ও অবহেলার কারণ}

\subsubsection*{কুরআনের সতর্কতা}

\begin{itemize}
\item \textbf{মাউন ১০৭:৪-৫:} \lat{“}ধিক সেই নামাজিদের জন্য, যারা নিজেদের নামাজ সম্পর্কে উদাসীন\lat{”} — যারা পড়ে কিন্তু অবহেলা করে; নামাজের সময়, রুকন বা মানের প্রতি উদাসীনতা ধিক্কারের কারণ।
\item \textbf{মারইয়াম ১৯:৫৯:} \lat{“}অতঃপর তাদের পরে এল অপদার্থ উত্তরসূরি, যারা নামাজ নষ্ট করল ও প্রবৃত্তির অনুসরণ করল; সুতরাং তারা অচিরেই ক্ষতির সম্মুখীন হবে।\lat{”} — নামাজ নষ্ট করার ভয়াবহ পরিণতি।
\end{itemize}
\subsubsection*{অবহেলার সাধারণ কারণ ও সমাধান}

\begin{table}[h]\centering\small
\begin{tabular}{ll}
কারণ & সমাধান \\
\hline
সময়ের অজ্ঞানতা/ভুলে যাওয়া & আযান-অ্যাপ, মোবাইল \textbf{\lat{reminder} (স্মরণক)}; নামাজের সময় ঘড়ি মানা \\
অলসতা (বিশেষত ফজর) & তাড়াতাড়ি ঘুমানো; ফজরের আগে ওঠার নিয়ম; জামাতের সঙ্গী \\
খুশুহীনতা — নামাজকে বোঝা ভাবা & অর্থ বুঝে পড়া; নামাজের মর্যাদার আয়াত-হাদিস পড়া \\
ব্যস্ততার অজুহাত & নামাজকে দিনের \textbf{\lat{priority} (অগ্রাধিকার)} বানানো — সময় সচেতনভাবে বরাদ্দ \\
পাপের কারণে দূরে থাকা & তাওবা করে ফিরে আসা — নামাজই পাপ থেকে বিরত রাখে (আনকাবুত ২৯:৪৫) \\
\hline
\end{tabular}
\end{table}

\medskip\hrule\medskip

\subsection*{৬. সালাতের গুণাবলি (\lat{Virtues}) — সহীহ দলিলসহ}

\begin{enumerate}
\item \textbf{জামাতের মর্যাদা:} \lat{“}জামাতে নামাজ একাকী নামাজের চেয়ে ২৫/২৭ গুণ মর্যাদাপূর্ণ\lat{”} (বুখারি ৬৪৫, মুসলিম ৬৫০); \lat{“}জামাতে ফজর পড়লে যেন সারা রাত ইবাদত করল\lat{”} (মুসলিম ৬৫৬)।
\item \textbf{প্রথম কাতার ও জামাতে যোগদান:} \lat{“}যদি মানুষ জানত আজান ও প্রথম কাতারে কী আছে...\lat{”} (বুখারি ৬১৫) — আযান ও প্রথম কাতারের প্রতি তাগিদ।
\item \textbf{আযানের জবাব ও দোয়া:} আযানের পরের দোয়ার বিনিময়ে সুপারিশের প্রতিশ্রুতি (বুখারি ৬১৪)।
\item \textbf{ফজরের দুই রাকাত সুন্নত:} দুনিয়ার চেয়ে উত্তম (মুসলিম ৭২৫)।
\item \textbf{ফেরেশতাদের দোয়া:} \lat{“}যতক্ষণ কেউ ওযুসহ মসজিদে বসে থাকে, ফেরেশতারা তার জন্য দোয়া করতে থাকে\lat{”} (বুখারি ৪৪৫, ৬৪৭; মুসলিম ৬৪৯)।
\item \textbf{পাঁচ নামাজের কাফফারা:} (মুসলিম ২৩৩) — দৈনন্দিন পাপ মোচনের নিয়মিত সুযোগ।
\item \textbf{তাহাজ্জুদ:} \lat{“}ফরজের পর সর্বোত্তম নামাজ রাতের নামাজ\lat{”} (মুসলিম ১১৬৩); নবীজির পা ফুলে যাওয়া পর্যন্ত দাঁড়ানো (বুখারি ৪৮৩৬, মুসলিম ২৮১৯)।
\end{enumerate}
\medskip\hrule\medskip

\subsection*{৭. সাহাবি ও সালাফের মনোভাব}

সাহাবারা নামাজকে জীবনের কেন্দ্রে রেখেছিলেন — এর প্রমাণ তাদের জীবনের অসংখ্য ঘটনা। এই সংকলনের `সাহাবিদের-ঘটনা/` ফোল্ডারে সূত্রসহ বিস্তারিত আছে; সংক্ষেপে:

\begin{itemize}
\item \textbf{নবীজি (সা.)} — তাহাজ্জুদে পা ফুলে যাওয়া; \lat{“}আমি কি কৃতজ্ঞ বান্দা হব না?\lat{”} (বুখারি ৪৮৩৬, ৬৪৭১; মুসলিম ২৮১৯)
\item \textbf{আবু বকর (রাঃ)} — নবীজির জীবদ্দশায় নামাজে ইমামতি (বুখারি ৬৮৩)
\item \textbf{উমার (রাঃ)} — ফজরের জামাতের ইমামতির সময় শাহাদাত (বুখারি ৩৭০০)
\item \textbf{হুযাইফা (রাঃ)} — নবীজির পেছনে দীর্ঘ রাতের নামাজের বর্ণনা (মুসলিম ৭৭২)
\item \textbf{ইবনে আব্বাস (রাঃ)} — খালার ঘরে নবীজির সাথে তাহাজ্জুদ (বুখারি ১১৭; মুসলিম ৭৬৩)
\end{itemize}
এই ঘটনাগুলোর সাধারণ শিক্ষা — সাহাবারা নামাজকে \textbf{সময় নষ্ট (\lat{waste} \lat{of} \lat{time})} ভাবতেন না; নামাজই ছিল তাদের দিনের কাঠামো, বিপদের আশ্রয় ও আল্লাহর সাথে কথোপকথন।

\medskip\hrule\medskip

\subsection*{৮. সারসংক্ষেপ (\lat{Summary})}

\begin{itemize}
\item নামাজ ইসলামের \textbf{স্তম্ভ (\lat{pillar})} — কুরআন ও সহীহ হাদিসে সর্বোচ্চ মর্যাদার ইবাদত।
\item নামাজ ঈমানের সাথে যুক্ত — ত্যাগকারীর জন্য কঠোর সতর্কতা; মতভেদ শাস্তির শ্রেণি নিয়ে, গুরুত্ব নিয়ে নয়।
\item নামাজ \textbf{হিসাবের (\lat{accountability})} প্রথম বিষয় — সময়মতো, পরিপূর্ণ ও খুশুসহ আদায় কাম্য।
\item নামাজের গুণাবলি অসংখ্য — জামাত, খুশু, সময়ের প্রতি যত্নে সওয়াব বহুগুণ।
\item সাহাবি ও সালাফের জীবনই নামাজের প্রকৃত মূল্যায়নের \textbf{উদাহরণ (\lat{example})}।
\end{itemize}
\begin{quote}
\textbf{মূল বার্তা:} নামাজকে ঘরে-মসজিদে, সুখে-দুঃখে, সফরে-বাসায় প্রতিষ্ঠা করা — এটিই মুমিন জীবনের প্রথম দায়িত্ব ও সর্বোচ্চ সাফল্যের চাবি।
\end{quote}

\medskip\hrule\medskip
\textit{সংশ্লিষ্ট ফাইল: `কুরআন/\lat{01}-কুরআনের-আয়াত.\lat{md}`, `হাদিস/\lat{02}-হাদিস-ও-আসার.\lat{md}`, `সাহাবিদের-ঘটনা/`}


\end{document}
